% =====================================================================
\chapter{Relaciones de equivalencia y de orden}
\label{tema:5}
% =====================================================================

Una relación binaria expresa que dos elementos están «relacionados» de cierta manera: ser iguales módulo~$n$, tener la misma longitud, preceder a otro. Este tema estudia las dos familias de relaciones más importantes, las de equivalencia, que clasifican, y las de orden, que ordenan. Sigue el capítulo~11 de Hammack~\cite{hammack}, las secciones~4.2 a~4.5 de Velleman~\cite{velleman} y las secciones~4.4 y~9.6 a~9.10 de~\cite{mcs}.

\section{Relaciones binarias}

\begin{definicion}[Relación]
Una \emph{relación} (binaria) en un conjunto $A$ es un subconjunto $R\subseteq A\times A$. Se escribe $a\,R\,b$ si $(a,b)\in R$.
\end{definicion}

\addtocounter{definicion}{1}%

Una relación en un conjunto finito puede representarse con un grafo dirigido, con un arco $a\to b$ por cada par $(a,b)\in R$, o con una matriz de ceros y unos cuya entrada $(a,b)$ vale~1 si $a\,R\,b$.

\begin{definicion}[Propiedades]\label{def:t5-prop}
Una relación $R$ en $A$ es
\begin{enumerate}
  \item \emph{reflexiva} si $a\,R\,a$ para todo $a\in A$;
  \item \emph{simétrica} si $a\,R\,b$ implica $b\,R\,a$;
  \item \emph{antisimétrica} si $a\,R\,b$ y $b\,R\,a$ implican $a=b$;
  \item \emph{transitiva} si $a\,R\,b$ y $b\,R\,c$ implican $a\,R\,c$.
\end{enumerate}
\end{definicion}

\begin{center}\small
\begin{tabular}{@{}lcccc@{}}
\toprule
Relación & Reflexiva & Simétrica & Antisimétrica & Transitiva \\
\midrule
$=$ en $\Z$ & sí & sí & sí & sí \\
$\le$ en $\Z$ & sí & no & sí & sí \\
$<$ en $\Z$ & no & no & sí & sí \\
$\mid$ en $\N$ & sí & no & sí & sí \\
$\equiv$ (mód.\ $n$) en $\Z$ & sí & sí & no & sí \\
$\neq$ en $\Z$ & no & sí & no & no \\
\bottomrule
\end{tabular}
\end{center}

Simetría y antisimetría \emph{no} son contrarias: la igualdad tiene ambas, y la relación $\{(1,2),(2,1),(1,3)\}$ en $\{1,2,3\}$ no tiene ninguna.

\section{Relaciones de equivalencia y particiones}

\begin{definicion}[Relación de equivalencia]
Una relación $\sim$ en $A$ es de \emph{equivalencia} si es reflexiva, simétrica y transitiva. La \emph{clase de equivalencia} de $a\in A$ es $[a]=\{x\in A: x\sim a\}$, y el \emph{conjunto cociente} es $A/{\sim}=\{[a]: a\in A\}$.
\end{definicion}

\begin{definicion}[Partición]
Una \emph{partición} de $A$ es una familia de subconjuntos de $A$ no vacíos, disjuntos dos a dos y cuya unión es $A$.
\end{definicion}

\begin{teorema}\label{teo:t5-particion}
Sea $\sim$ una relación de equivalencia en $A$. Entonces:
\begin{enumerate}
  \item $a\in[a]$ para todo $a$;
  \item $[a]=[b]$ si y solo si $a\sim b$;
  \item dos clases son iguales o disjuntas.
\end{enumerate}
En consecuencia, $A/{\sim}$ es una partición de $A$. Recíprocamente, toda partición $\mathcal{C}$ de $A$ define una relación de equivalencia, $a\sim b$ si $a$ y $b$ están en el mismo bloque de $\mathcal{C}$, cuyas clases son los bloques.
\end{teorema}

\addtocounter{definicion}{1}%

\begin{proposicion}[Núcleo de una función]\label{prop:t5-nucleo}
Sea $f\colon A\to B$. La relación $a\sim_f a'$ si $f(a)=f(a')$ es de equivalencia, y sus clases son las fibras no vacías $f^{-1}(\{b\})$, $b\in f(A)$.
\end{proposicion}

La demostración es el \cref{eje:5-nucleo}. Muchas relaciones de equivalencia son de este tipo: «tener la misma longitud» es el núcleo de $\ell\colon\Sigma^{*}\to\N$, y la congruencia módulo $n$ es el núcleo de la función «resto de dividir entre $n$».

\section{Relaciones de orden}

\begin{definicion}[Orden]
Una relación $\preceq$ en $A$ es un \emph{orden parcial} si es reflexiva, antisimétrica y transitiva; el par $(A,\preceq)$ es un \emph{conjunto ordenado}. Es un \emph{orden total} si, además, dos elementos cualesquiera son comparables: $a\preceq b$ o $b\preceq a$. Un elemento $m$ es \emph{minimal} si no existe $x\neq m$ con $x\preceq m$, y \emph{mínimo} si $m\preceq x$ para todo $x$; análogamente se definen maximal y máximo.
\end{definicion}

\addtocounter{definicion}{1}%

Un conjunto ordenado finito se representa con su \emph{diagrama de Hasse}: se dibuja un segmento ascendente de $x$ a $y$ cuando $x\prec y$ y no hay ningún $z$ con $x\prec z\prec y$; la reflexividad y la transitividad quedan implícitas (\cref{fig:t5-hasse}).

\addtocounter{figure}{1}%

\addtocounter{definicion}{1}%

\addtocounter{definicion}{1}%

